\section{Theoretical Foundation}

\label{sec:theory}

We ground our approach in Hamiltonian mechanics, providing a principled framework for understanding the trade-offs in context-based learning.

\subsection{The Context-Inference Hamiltonian}

\begin{definition}[Policy Mass]
The \textbf{policy mass} $\Psi$ quantifies the semantic richness of the context provided to an LLM:
\begin{equation}
\Psi = \sum_{e \in \mathcal{E}} w_e \cdot \text{relevance}(e, q)
\end{equation}
where $\mathcal{E}$ is the experience library, $w_e$ is the weight (confidence) of experience $e$, and $\text{relevance}(e, q)$ measures applicability to query $q$.
\end{definition}

\begin{definition}[Inference Complexity]
The \textbf{inference complexity} $\Theta$ quantifies the computational cost and uncertainty:
\begin{equation}
\Theta = H[\pi_\theta(\cdot | q, \mathcal{E})] + \lambda \cdot |\mathcal{E}|
\end{equation}
where $H[\pi_\theta]$ is the policy entropy (model uncertainty) and $|\mathcal{E}|$ is the context length cost with regularization $\lambda$.
\end{definition}

\begin{theorem}[Hamiltonian Invariant]
For a sufficiently capable foundation model $\pi_\theta$ with frozen parameters $\theta$, there exists a conserved quantity $\kappa$ such that:
\begin{equation}
\Psi \cdot \Theta = \kappa
\end{equation}
This conservation law ensures that increasing policy mass (richer context) decreases inference complexity (lower uncertainty), maintaining system equilibrium.
\end{theorem}

\begin{proof}[Proof Sketch]
Consider the expected performance of the LLM as a function of both context and model uncertainty:
\begin{equation}
\mathbb{E}_{q \sim \mathcal{D}}[R(q, \pi_\theta(\cdot | q, \mathcal{E}))]
\end{equation}

For a fixed level of expected performance (target reward), the information-theoretic bound from Rate-Distortion theory \cite{cover1999elements} implies:
\begin{equation}
I(Q; O | \mathcal{E}) + H[O | Q, \mathcal{E}] \geq C
\end{equation}
where $I(Q; O | \mathcal{E})$ is mutual information between queries and outputs given context, $H[O | Q, \mathcal{E}]$ is conditional entropy, and $C$ is a constant determined by the target performance level.

The mutual information term $I(Q; O | \mathcal{E})$ increases with policy mass $\Psi$ (richer context provides more information), while the conditional entropy $H[O | Q, \mathcal{E}]$ corresponds to our inference complexity $\Theta$. The conservation law $\Psi \cdot \Theta = \kappa$ emerges as a geometric mean constraint that maintains this information-theoretic balance.
\end{proof}

\subsection{Policy Manifold Geometry}

Traditional fine-tuning navigates the \textbf{parameter space} $\mathcal{M}_\theta$, a high-dimensional manifold where each point represents a different set of model weights. Training-Free GRPO instead navigates the \textbf{context space} $\mathcal{M}_\mathcal{E}$, where each point represents a different experience library.

\begin{proposition}[Context Space Optimization]
For a frozen foundation model $\pi_\theta$ with sufficient capability, the optimal experience library $\mathcal{E}^*$ can be found by gradient-free optimization in context space:
\begin{equation}
\mathcal{E}^* = \arg\max_{\mathcal{E}} \mathbb{E}_{q \sim \mathcal{D}}[R(q, \pi_\theta(\cdot | q, \mathcal{E}))] - \lambda \cdot |\mathcal{E}|
\end{equation}
This optimization can be performed through discrete updates (Add/Modify/Delete operations) guided by semantic advantages.
\end{proposition}

The key insight: while $\mathcal{M}_\theta$ has billions of dimensions (model parameters), $\mathcal{M}_\mathcal{E}$ has only hundreds (experiences), making optimization vastly more efficient. Furthermore, $\mathcal{M}_\mathcal{E}$ is \textit{discrete and interpretable}—each point corresponds to a human-readable set of strategic insights.

\subsection{Semantic Advantage as Gradient Analog}

In traditional policy gradient methods, we compute:
\begin{equation}
\nabla_\theta J(\theta) = \mathbb{E}_{\tau \sim \pi_\theta}[A^\pi(s, a) \nabla_\theta \log \pi_\theta(a | s)]
\end{equation}
where $A^\pi(s, a)$ is the advantage function.

In Training-Free GRPO, we instead compute \textbf{semantic advantages}:
\begin{equation}
\mathcal{A}_{\text{sem}}(\tau_i, \tau_j) = \text{LLM}_{\text{introspect}}(\tau_i, \tau_j, r_i, r_j)
\end{equation}
which returns a natural language description of \textit{why} trajectory $\tau_i$ outperformed $\tau_j$.

This semantic advantage serves as an analog to gradients:
\begin{itemize}
    \item \textbf{Gradient}: Direction in parameter space to improve policy
    \item \textbf{Semantic Advantage}: Insight about context to add/modify for improvement
\end{itemize}

The "update" is then performed discretely in context space:
\begin{equation}
\mathcal{E}_{t+1} = \text{Update}(\mathcal{E}_t, \mathcal{A}_{\text{sem}})
\end{equation}
where Update is a discrete operation (Add/Modify/Delete) rather than continuous gradient descent.

\subsection{Convergence Analysis}

\begin{theorem}[Context Space Convergence]
Under mild regularity conditions (experience library size $|\mathcal{E}| \leq M$, bounded relevance scores, sufficiently capable foundation model), the Training-Free GRPO algorithm converges to a local optimum in context space in $\mathcal{O}(M \log M)$ iterations.
\end{theorem}

\begin{proof}[Proof Sketch]
Each iteration performs group-relative comparisons across $G$ rollouts, extracting semantic advantages that strictly improve the experience library (measured by expected reward). Since the context space is discrete and finite (bounded by $M$ experiences, each with finite vocabulary), and each update increases expected performance, the algorithm must converge to a fixed point where no further advantageous experiences can be extracted. The $\log M$ factor arises from the tree structure of embedding-based retrieval for relevant experiences.
\end{proof}

Importantly, this convergence requires \textit{no gradient computation}, making it applicable to black-box API models where parameter access is unavailable.

